% !TeX root = ../main.tex

\chapter{Introduction}\label{chapter:introduction}

\section{Motivation}\label{sec:motivation}

Contract Lifecycle Management (CLM) is the connective tissue of modern enterprise
commerce. Every commercial relationship, employment arrangement, vendor engagement,
and regulatory commitment is encoded in a contract, and the operational performance
of a firm depends on its ability to draft, review, execute, monitor, and renew these
documents at scale. Industry surveys conducted by World Commerce \& Contracting
indicate that organisations lose value equivalent to roughly nine percent of annual
revenue as a consequence of poor contracting practices, with the figure rising to
fifteen percent or more in industries with high contractual
complexity~\cite{worldcc2025benchmark}. The same survey reports that contract-related
data is, on average, dispersed across twenty-four distinct systems, and that
eighty-three percent of executives consider their contracts too rigid to adapt to
changing business conditions~\cite{worldcc2025benchmark}. These figures position CLM
not as a clerical concern but as a material driver of enterprise financial performance.

The cost structure of manual contract review compounds this strategic exposure. Legal
professionals spend a substantial proportion of their billable time on tasks that are,
in principle, automatable: locating clauses, comparing them against playbooks, flagging
deviations, and tracking obligations. Goldman Sachs research estimates that
approximately forty-four percent of legal-task time in the United States is technically
exposed to automation by generative artificial intelligence, the second-highest exposure
rating across surveyed professions~\cite{goldmansachs2023ai}. Vendor-led empirical
comparisons such as the Onit-funded study by Martin et al.\ report that frontier large
language models match or exceed junior lawyers and legal-process outsourcers on
issue-identification F-scores while reducing per-document review time by roughly two
orders of magnitude~\cite{martin2024bettercallgpt}. Although such vendor evidence must
be read with appropriate skepticism, the directional finding --- that LLMs are
competitive on standardised review tasks --- is consistent across independent academic
benchmarks~\cite{hendrycks2021cuad,liu2025contracteval}.

The appeal of LLM automation, however, is not the same as the demonstration of its
reliability under enterprise conditions. State-of-the-art models still hallucinate
clauses that are not present, omit clauses that are, exhibit pronounced sensitivity to
prompt phrasing~\cite{chatterjee2024posix}, and fail in ways that are difficult to
audit. Liu et al.\ document a \emph{laziness} failure mode in which open-source models
default to ``no relevant clause'' responses even when the relevant text is plainly
present in the contract~\cite{liu2025contracteval}. Single-prompt zero-shot
deployments, which dominate current commercial CLM offerings, give product teams
little ability to localise these failures, attribute them to specific model behaviours,
or remediate them without retraining or wholesale model substitution.

In response, both the research community and the agentic-AI vendor ecosystem have
advocated multi-agent architectures as the next paradigm for legal document automation.
The argument is that decomposing contract analysis into specialised cooperating
agents --- a classifier, an extractor, a reviewer, a router, a verifier --- ought to
improve accuracy through specialisation, enable explainability through inspectable
inter-agent traces, and accommodate enterprise governance requirements such as those
imposed by Article~13 of the European Union's Artificial Intelligence
Act~\cite{euaiact2024} and Article~22 of the General Data Protection
Regulation~\cite{gdpr2016}. Multi-agent frameworks such as
AutoGen~\cite{wu2023autogen}, MetaGPT~\cite{hong2024metagpt},
CAMEL~\cite{li2023camel}, and Mixture-of-Agents~\cite{wang2024moa} have made these
architectures cheap to construct, and orchestration libraries such as
LangGraph~\cite{langgraph2024} have lowered the engineering threshold further. Whether
these architectures actually \emph{work} for contract analysis at production scale is,
however, an open empirical question --- and it is the question this thesis takes up.

\section{Research Gap}\label{sec:research-gap}

Three observations motivate the research gap that this thesis addresses.

First, the empirical evidence base for which multi-agent patterns improve contract
analysis is thin. Recent legal multi-agent systems such as
ChatLaw~\cite{cui2024chatlaw}, PAKTON~\cite{raptopoulos2025pakton}, and
L-MARS~\cite{wang2025lmars} are valuable system-design contributions, but each
evaluates a single architecture on a single benchmark, typically without ablating the
architectural choice itself. There is, to our knowledge, no published cross-architecture,
cross-model-family comparison on a public contract-analysis benchmark that isolates the
contribution of the multi-agent design from the contribution of the underlying LLM.

Second, recent work on multi-agent systems outside the legal domain has begun to
question the universality of multi-agent gains. Gao et al.\ argue that the benefits of
multi-agent systems over single-agent systems diminish as base-model capability
increases, because frontier models with stronger long-context reasoning, memory, and
tool-use absorb the workload that previously required
decomposition~\cite{gao2025single}. Cemri et al.\ construct a taxonomy of fourteen
failure modes specific to multi-agent LLM systems, grouped into specification,
inter-agent misalignment, and verification-and-termination categories, and report that
simple interventions yield only modest gains~\cite{cemri2025why}. Bogavelli et al., in
their AgentArch benchmark, find that no agent architecture is universally optimal across
model families and that smaller models can outperform larger models when matched to an
appropriate architecture~\cite{bogavelli2025agentarch}. These results suggest that
whether multi-agent helps in CLM depends on conditions that have not yet been mapped.

Third, prior contract-analysis evaluations have used sampling and metric choices that
may obscure the dominant failure mode. The CUAD benchmark exhibits a heavily imbalanced
class distribution at the contract-clause level, but several published evaluations
resample to a positive-heavy ratio, which produces optimistic precision estimates and
hides over-extraction failures~\cite{liu2025contracteval}. Sampling design and metric
weighting interact in ways that change the ranking of architectures, and these
interactions have not been systematically interrogated.

Taken together, these observations indicate a gap that is simultaneously empirical,
mechanistic, and methodological: the field lacks a rigorous comparison of multi-agent
architectural patterns for clause-level contract analysis under realistic class
distributions, with statistical hypothesis tests, mechanistic diagnosis of failure
modes, and explicit treatment of when multi-agent designs help versus when they do not.

\section{Research Questions}\label{sec:research-questions}

This thesis investigates the following main research question:

\begin{quote}
\emph{Under what conditions do multi-agent architectures improve contract analysis
accuracy, efficiency, and explainability in enterprise CLM systems, and which
architectural patterns are most effective?}
\end{quote}

Three sub-questions decompose this main question into empirically tractable parts.

\textbf{Sub-RQ1.} Which multi-agent architectural patterns --- pipeline decomposition,
specialist routing, or consolidated prompting --- are most effective for clause-level
contract analysis?

\textbf{Sub-RQ2.} How do multi-agent approaches compare to single-agent baselines in
clause extraction accuracy and cost-efficiency across diverse clause categories and
contract lengths?

\textbf{Sub-RQ3.} What architectural decisions enable auditable reasoning traces in
multi-agent contract analysis, and what design tradeoffs do they introduce?

The first sub-question targets architectural ranking conditional on model family. The
second sub-question incorporates economic and class-distribution considerations that
govern enterprise deployability. The third sub-question treats explainability as a
system-design property rather than as a metric to be optimised, in line with the
regulatory framing imposed by Articles~13 and~86 of the EU AI Act~\cite{euaiact2024}.

\section{Hypotheses}\label{sec:hypotheses}

The empirical core of the thesis tests three hypotheses. A fourth hypothesis on
traceability, present in the original thesis proposal, has been reframed: rather than
test whether multi-agent designs produce more inspectable traces --- a property that is
largely determined by the orchestration layer rather than the architecture per se --- we
treat traceability as a system property of the implementation and document the design
choices that secure it (see Sub-RQ3).

\textbf{H1.} \emph{At least one multi-agent architecture outperforms the strongest
single-agent baseline on $F_2$.} The rationale is that decomposition should reduce
per-step prompt length and isolate the classification decision from the extraction
decision, which has been argued to improve recall in imbalanced retrieval
tasks~\cite{liu2025contracteval}. We test H1 with a paired McNemar's
test~\cite{mcnemar1947note} between the best multi-agent configuration and the best
single-agent configuration per model, with effect sizes reported.

\textbf{H2.} \emph{Multi-agent architectures reduce the rare-common $F_2$ performance
gap compared to single-agent baselines.} CUAD's clause categories vary substantially in
prevalence, and architectures that route to category-specialised prompts may benefit
rare categories disproportionately. We test H2 by comparing the rare-versus-common
$F_2$ gap across architectures using a paired bootstrap confidence interval on the gap
difference.

\textbf{H3.} \emph{Specialist architecture differs significantly from consolidated
specialist prompts.} If routing-plus-specialists outperforms a single prompt that
contains all specialist guidance, the difference must be attributable to the routing or
to the per-specialist context isolation. We test H3 with a two-sided McNemar's test
contrasting the LLM-routed multi-agent configuration against a single-agent ablation
that contains the union of specialist instructions, and additionally contrast the
pipeline-decomposition configuration against the consolidated prompt.

The reader is alerted that the empirical results, reported in Chapter~5, are mixed with
respect to these hypotheses, and that the thesis treats the negative or null findings as
substantive contributions rather than as failures of the design.

\section{Contributions}\label{sec:contributions}

This thesis advances five contributions, which we summarise here and develop in
subsequent chapters.

The first contribution is methodological. We introduce a natural class-distribution
sampling correction for CUAD-based evaluation, in which the positive-to-negative ratio
of clause-presence labels is held at approximately thirty to seventy, matching CUAD's
underlying distribution rather than the positive-heavy resampling used in some prior
work. This correction exposes \emph{over-extraction} as the dominant failure mode of
zero-shot LLM clause extraction on long contracts, a mode that is suppressed by
positive-heavy sampling. The implication is that prior reported precision figures for
LLM contract review may overstate enterprise-relevant performance.

The second contribution is empirical. We report what is, to our knowledge, the first
cross-family cross-architecture comparison on CUAD that simultaneously varies
architecture and frontier-model identity. We evaluate five architectural configurations
spanning zero-shot, chain-of-thought, pipeline decomposition, LLM-routed specialists,
and a consolidated single-agent ablation, against eleven frontier models from the
Anthropic, Google, and OpenAI families, on a fixed sample of approximately 2{,}999
instances drawn from eighty-one CUAD contracts longer than 100{,}000 characters. The
resulting matrix yields 163{,}644 sample-level predictions analysed with paired
statistical tests.

The third contribution is mechanistic. We identify and name a failure mode specific to
pipeline-decomposition multi-agent architectures: the \emph{classification-gate
rejection cascade}. When a first-stage classifier overrejects on borderline cases,
downstream extractors are never invoked, recall collapses, and the failure is invisible
from the perspective of the extractor's own metrics. We document this cascade
quantitatively and trace it to a small number of model-family interactions.

The fourth contribution is an architectural critique of the LLM-routed specialist
pattern as it appears in current legal multi-agent systems. We observe that on the
contract-analysis task, the LLM-router achieves near-perfect routing accuracy (one
hundred percent on ten of eleven models in our evaluation), the routed configuration is
statistically indistinguishable from a consolidated single-agent prompt on $F_2$, and
the routed configuration produces less clean span outputs. The implication, which we
develop in Chapter~6, is that the LLM-routed specialist pattern, in the form most often
advocated, can be replaced with a deterministic dictionary lookup without measurable
accuracy loss.

The fifth contribution is practitioner-oriented: a six-gap analysis connecting CUAD
benchmark performance to enterprise CLM deployment reality. We mention this contribution
here without elaboration, as it is developed in Chapter~6.

We emphasise here, to set reader expectations, that the empirical findings of the thesis
are not uniformly positive for multi-agent designs. No multi-agent configuration we
evaluate consistently outperforms the strongest single-agent baseline on $F_2$, and the
mean $F_2$ across eleven models for all five configurations falls within a 1.1-point
band. Architectural ranking varies systematically by model family. The thesis frames
these results not as an indictment of multi-agent CLM but as a delineation of the
conditions under which multi-agent designs help, hurt, or are inert --- which is the
question Sub-RQ1 explicitly asks.

\section{Thesis Outline}\label{sec:thesis-outline}

The remainder of the thesis is organised as follows. Chapter~2 reviews the literature
on contract analysis, the evolution from rule-based to neural to LLM-based approaches,
multi-agent systems for legal AI, the conditions under which multi-agent designs help or
hurt, and evaluation methodology in legal NLP, and positions the present work against
this body of evidence. Chapter~3 specifies the methodology, including the five
architectural configurations evaluated, the eleven frontier models surveyed, the natural
class-distribution sampling design, the metric set, and the statistical testing
protocol. Chapter~4 describes the implementation, including the LangGraph-based
orchestration~\cite{langgraph2024}, the Langfuse-based observability
layer~\cite{langfuse2024}, prompt construction, and reproducibility considerations.
Chapter~5 reports the empirical results, the per-hypothesis tests, and the mechanistic
diagnosis of the classification-gate rejection cascade. Chapter~6 develops the
architectural critique of the LLM-routed specialist pattern, the practitioner-oriented
six-gap analysis, and a discussion of threats to validity. Chapter~7 concludes with
limitations, implications for enterprise CLM SaaS design, and directions for future
work.
